Современный корпоративный маркетинговый сайт выполняет не только представительскую,
но и операционную функцию: он используется для публикации страниц под маркетинговые
кампании, размещения статей и кейсов, публикации вакансий, сбора откликов и поддержки
поисковой видимости компании в цифровой среде. При этом ценность такой платформы
определяется не только качеством визуального интерфейса, но и скоростью обновления
контента, предсказуемостью публикации и возможностью передавать редакторские задачи
контент-менеджеру без постоянного участия разработчика.

Для маркетинговых сайтов типична ситуация, при которой публичная витрина развивается
быстрее, чем механизмы управления содержимым. В результате часть контента и интерфейсных
текстов остается зашитой во frontend-коде, публикация изменений требует ручного участия
разработчика, а сценарии локализации, предварительного просмотра и управления
\texttt{SEO}-метаданными оказываются фрагментированными. Такое состояние затрудняет
поддержку сайта как контентной системы и снижает эффективность взаимодействия между
разработкой и маркетинговой функцией компании.

\frontsection{Актуальность темы}

Актуальность темы обусловлена тем, что корпоративные маркетинговые сайты требуют
регулярного обновления контента, поддержки нескольких языков, управления структурой
страниц, публикации материалов с учетом \texttt{SEO}-требований и воспроизводимого
сценария вывода изменений на публичную витрину. При монолитном или частично
кодо-центричном подходе к разработке сайта эти задачи решаются разрозненно: одни данные
хранятся в \texttt{CMS}, другие остаются во frontend-приложении, а публикационный процесс
зависит от ручных действий и знания внутренней структуры проекта. Такая постановка
подтверждается как исследованиями роли сайта в маркетинговых коммуникациях, так и
прикладными работами по требованиям к \texttt{CMS} в retail-контуре~\cite{torgunakova2023marketing-website,voroniuk2020retail-cms-needs}.

Для рассматриваемой предметной области такое противоречие проявляется между потребностью
в быстром редакторском управлении маркетинговым контентом и ограничениями архитектуры,
в которой управление содержанием недостаточно отделено от слоя отображения. В этих
условиях актуальной становится разработка \texttt{CMS-first} системы, где \texttt{Strapi}
используется как ядро управления контентом, а \texttt{Astro} --- как публичная витрина,
ориентированная на производительный рендеринг и контролируемую публикацию.

Дополнительную значимость для данной предметной области имеет возможность предварительной
сборки публичных страниц frontend-слоя. Для маркетингового сайта, в котором значительная
часть материалов представляет собой контентные страницы с относительно редким изменением
между публикациями, такой подход позволяет отдавать заранее подготовленное представление
страницы, уменьшать зависимость витрины от runtime-логики \texttt{CMS} и повышать
предсказуемость публикационного контура. В работе этот эффект рассматривается как
архитектурное следствие связки \texttt{Strapi + Astro} и схемы
\texttt{webhook -> rebuild}, а не как абстрактное свойство любой \texttt{CMS}.

Практическая значимость такой постановки усиливается тем, что работа выполняется не как
абстрактное исследование технологий, а как разработка проектной реализации в
монорепозитории \texttt{Nx + pnpm}, включающем backend на
\texttt{Strapi 5} и frontend на \texttt{Astro 6}. Это позволяет рассматривать тему ВКР
как инженерную задачу построения целостной модели управления страницами, статьями,
проектами, вакансиями и откликами с учетом локализации, режима предварительного
просмотра, \texttt{SEO}, схемы \texttt{webhook -> rebuild} и контура развертывания.

\frontsection{Объект и предмет исследования}

Объектом исследования являются процессы управления контентом корпоративного маркетингового
сайта.

Предметом исследования являются архитектурные и программные решения для построения
\texttt{CMS-first} веб-платформы на базе \texttt{Strapi} и \texttt{Astro}, обеспечивающей
редакторское управление страницами, статьями, проектами, вакансиями и откликами без
непосредственного вмешательства в исходный код публичной витрины.

\frontsection{Цель и задачи работы}

Целью работы является разработка системы управления контентом для корпоративного
маркетингового сайта на базе \texttt{Strapi} и \texttt{Astro}, обеспечивающей управление
страницами, статьями, кейсами и вакансиями, поддержку локалей \texttt{ru} и \texttt{en},
предварительного просмотра, \texttt{SEO}-настроек и воспроизводимого сценария публикации
без постоянного участия разработчика.

Для достижения поставленной цели необходимо решить следующие задачи:
\begin{enumerate}
  \item проанализировать особенности маркетинговых сайтов, подходы к управлению
  контентом и обосновать выбор \texttt{CMS-first} архитектуры на базе \texttt{Strapi} и
  \texttt{Astro};
  \item сформировать функциональные и нефункциональные требования к системе, а также
  спроектировать ее архитектуру и модель данных для ключевых контентных сущностей;
  \item реализовать \texttt{CMS-first} контур управления страницами и ключевым контентом,
  включая коллекцию \texttt{pages}, конструктор страниц на основе \texttt{Dynamic Zone} и
  вынос основных интерфейсных текстов в \texttt{Strapi};
  \item реализовать ключевые возможности публичной витрины, включая поддержку локалей
  \texttt{ru/en}, режим предварительного просмотра, управление \texttt{SEO}-полями,
  \texttt{Open Graph} и генерацию \texttt{sitemap};
  \item организовать эксплуатационный контур системы, включающий роли и права редакторов,
  защиту публичных сценариев, публикацию через схему \texttt{webhook -> rebuild}, а также
  deployment frontend и \texttt{CMS} как отдельных \texttt{Docker}-приложений под
  управлением \texttt{Dokploy};
  \item провести проверку ключевых пользовательских и редакторских сценариев и оценить
  полученный результат по функциональным, локализационным, \texttt{SEO}- и
  эксплуатационным характеристикам.
\end{enumerate}

\frontsection{Методы исследования и разработки}

В работе используются методы системного и сравнительного анализа, применяемые для
исследования особенностей маркетинговых сайтов и выбора архитектурного подхода к
управлению контентом. Для проектирования системы используются методы структурного анализа
предметной области, моделирования контентных сущностей и их связей, а также
архитектурного проектирования распределенной веб-платформы с разделением контентного и
витринного слоев.

В практической части применяются методы компонентной разработки frontend-приложений,
интеграционной настройки \texttt{Strapi} и \texttt{Astro}, формирования типизированного
\texttt{API}-контракта на основе \texttt{OpenAPI}, конфигурирования deployment-процесса и
функционального тестирования сценариев публикации, локализации и работы публичных форм.
Дополнительно используются приемы инженерной валидации, основанные на сопоставлении
поставленных требований с фактическим состоянием кода, конфигурации и проверяемых
сценариев системы.

\frontsection{Теоретическая значимость}

Теоретическая значимость работы заключается в систематизации подхода к построению
\texttt{CMS-first} архитектуры для маркетингового сайта, в которой управление контентом,
локализация, публикация, поисковая оптимизация и публичный рендеринг рассматриваются как
взаимосвязанные части одной инженерной системы. Работа уточняет принципы разграничения
ответственности между \texttt{CMS}-ядром и frontend-витриной, а также показывает, каким
образом контентная модель, сценарии публикации и deployment-конфигурация должны
проектироваться согласованно, а не как набор изолированных доработок.

Дополнительная теоретическая ценность состоит в формализации требований к предметной
области корпоративного маркетингового сайта, где ключевое значение имеют не только
отображение контента, но и редакторские сценарии: создание страниц из блоков,
мультиязычность, предварительный просмотр, управление метаданными и воспроизводимая
публикация на публичную витрину.

\frontsection{Практическая значимость}

Практическая значимость работы определяется тем, что ее результат ориентирован на
непосредственное использование в реальном проекте. Разработанная система обеспечивает
перенос ключевого контента и редакторских операций в \texttt{Strapi}, снижает
зависимость маркетинговой команды от разработчика при типовых изменениях и обеспечивает
единый путь публикации материалов на публичный сайт.

Для контент-менеджера и маркетолога практическая ценность решения состоит в возможности
управлять страницами, статьями, проектами, вакансиями и связанными метаданными через
специализированный административный интерфейс, а не через редактирование исходного кода.
Для проекта в целом практическая значимость выражается в централизованной модели данных,
поддержке локалей \texttt{ru/en}, подготовке сценария предварительного просмотра,
организации \texttt{SEO}-данных и воспроизводимого контура развертывания, в котором
frontend и \texttt{CMS} развертываются как отдельные \texttt{Docker}-приложения под
управлением \texttt{Dokploy}. Разработанная в рамках ВКР система подтверждается кодом, конфигурацией и
воспроизводимыми тестовыми сценариями, поэтому ее практическая значимость проявляется как
реализованный инженерный результат.

Дополнительно практическая ценность выбранной архитектуры связана с тем, что публичная
витрина может собираться заранее на стороне frontend-слоя. Для контентных страниц это
создает предпосылки для более предсказуемой отдачи результата и снижает необходимость
формировать представление страницы при каждом обращении пользователя. Данный вывод в
работе трактуется как архитектурное преимущество выбранной схемы, а не как заранее
измеренное количественное утверждение о производительности.

\frontsection{Краткое описание результата работы}

Кратким результатом работы является разработка архитектурного и программного решения для
перехода корпоративного маркетингового сайта к \texttt{CMS-first} модели, в которой
\texttt{Strapi} выступает центром управления контентом, а \texttt{Astro} выполняет роль
публичной витрины. Целевая система охватывает управление страницами, статьями, проектами,
вакансиями и откликами, сборку страниц на основе \texttt{Dynamic Zone}, поддержку локалей
\texttt{ru/en}, режим предварительного просмотра, \texttt{SEO}, \texttt{Open Graph},
\texttt{sitemap}, роли и права, а также публикационный контур
\texttt{webhook -> rebuild}.

В рамках работы сформирован монорепозиторий с самостоятельными приложениями
\texttt{Strapi} и \texttt{Astro}, реализованными модулями статей, проектов, вакансий,
откликов, маркетинговых страниц и редакторских сценариев публикации. В тексте диплома
раскрываются архитектура системы, состав реализованных модулей, границы принятого объема и
подтверждение итогового результата архитектурой, конфигурацией и тестированием.

\frontsection{Структура работы}

Работа состоит из введения, двух глав, заключения, списка сокращений, словаря терминов,
списка использованных источников и приложений. В первой главе рассматриваются особенности
маркетинговых сайтов как объекта автоматизации, сравниваются традиционные и
\texttt{headless} подходы к управлению контентом, обосновывается выбор \texttt{Strapi} и
\texttt{Astro}, а также формируются требования к системе. Во второй главе описываются
архитектура решения, модель данных, сценарии публикации, подход к локализации, deployment
pipeline, вопросы безопасности и результаты тестирования. В заключении подводятся итоги
работы, формулируются достигнутые результаты, ограничения и направления дальнейшего
развития системы.
